\chapter{Conclusões}
%e Recomendações
\label{cap:05}

De forma geral, observou-se que tecnologias com maior nível de abstração, como \textit{Godot} e \textit{Pygame}, proporcionaram maior rapidez de desenvolvimento e menor complexidade inicial. Em contrapartida, tecnologias mais próximas do baixo nível, como \textit{Raylib} com \textit{C++}, exigiram maior conhecimento técnico, mas ofereceram maior controle no desenvolvimento.

É importante salientar que na escolha da tecnologia deve levar em consideração o objetivo do projeto, o nível de experiência do desenvolvedor e os requisitos específicos do jogo a ser desenvolvido. 
Para projetos mais simples ou para iniciantes, tecnologias como \textit{Pygame} ou \textit{Godot} podem ser mais adequadas. Já para projetos que demandam maior controle e performance, tecnologias como \textit{Raylib} podem ser mais indicadas, desde que o desenvolvedor possua o conhecimento necessário para lidar com a complexidade envolvida.

A experiência de desenvolver múltiplas versões do mesmo jogo permitiu compreender como diferentes tecnologias influenciam diretamente na organização do código, na complexidade do desenvolvimento e na implementação de mecânicas de jogos digitais.
